\SourceLeaf{216}$cbe = cb'b$. Thus $be$ will intersect 2 at $e'$. If one now makes $b'g = be'$, then $cg$ and $cd'$ will make equal angles with $cb'$. Which is absurd.

5. If 1 is parallel to 2 and 1 to 3, then 2 is also parallel to 3.

6. What corresponding points on two parallels are.

7. Equidistance of corresponding points.

8. The point on a third parallel corresponds to corresponding points on the first two.

9. \textls{Trope} is the l[ine formed by corresponding points when one considers all parallels to a straight line.]

\Needspace{19\baselineskip}
\bigskip
\begin{center}
\rule{2.5cm}{0.4pt}

\textsc{Remark.}
\end{center}

In the letter to \textsc{Schumacher} of 17 May 1831 (p. 213 of this volume), \textsc{Gauss} says that he had never previously written down anything from his reflections on the theory of parallel lines, which were already about 40 years old, but that he had begun to write down some of them a few weeks earlier. In the letter to \textsc{Bolyai} of 6 March 1832, \textsc{Gauss} also mentions such notes. One will therefore hardly be mistaken in assuming that the foregoing undated notes, which are on separate slips, date from 1831.

The figures in note [1, and the first figure in note [3], were added to the text

\hfill \textsc{Stäckel}.

